\chapter{Information Structure}\label{ch:infstr}
\label{sec:topic-comment}% cross-reference compatibility: retained for \ref{sec:topic-comment}

The analysis of Iridian information structure developed in this chapter
presupposes the clause spine and theoretical framework introduced in
§\,\ref{sec:grammatical-architecture} of Chapter~\ref{ch:overview} and the
background discussion in §\,\ref{sec:th-background} of
Chapter~\ref{ch:simple}. The orientation given there --- that the absolutive
pivot is strongly privileged as the default discourse centre, and that
topic-comment organisation is the primary mode of the Iridian clause --- is
taken as the point of departure. The present chapter provides the full
analytical account.


\section{The topic: conceptual and formal foundations}\label{sec:tc-foundations}

\subsection{The topic as the logical subject of predication}
\label{sec:tc-logical-subject}

The Iridian clause is organised around a \term{topic-comment bipartition}, and
the account of that bipartition must be precise enough to do real theoretical
work. The informal gloss --- topic is `what the sentence is about' --- is
insufficient, because it conflates two analytically separate relations: the
grammatical relation of subject, which is defined by morphosyntactic properties
such as case, agreement, and pivot behaviour, and the discourse relation of
topichood, which is defined by the semantic relation between a constituent and
the predication made about it. In the framework adopted here, following Kiss
(1995), the \term{topic} is the \term{logical subject of predication}: the
constituent with respect to which the rest of the clause --- the \term{comment}
--- constitutes a semantic predicate, expressing a property that is claimed to
hold of that individual. This is a discourse-semantic role, not a
morphosyntactic one, and it is independent of grammatical subjecthood in the
standard sense.

The independence of topichood from subjecthood is a theoretical necessity in
Iridian because the grammar does not recognise a language-wide grammatical
subject as its organising primitive. The stable privilege belongs instead to the
\term{absolutive pivot}: the direct-case core argument licensed by the alignment
system. In nominative-domain clauses --- the imperfective aspects --- the
absolutive pivot is typically the agent or sole participant. In ergative-domain
clauses --- the perfective and retrospective --- it is the patient. The same
argument position, determined by alignment, is the discourse centre in both
domains; but the semantic role of that centre changes with aspect domain.

This matters for topichood in the following way. If `topic' were simply a
synonym for `grammatical subject', the Iridian facts would require two different
subjects: an agentive subject in nominative-domain clauses and a patient subject
in ergative-domain clauses. That is not a coherent analysis. Instead, the
grammar treats \term{topic} as a discourse role assigned to a structurally
identified position, Spec,\synhead{TopP}, and the absolutive pivot as the
default occupant of that position across aspect domains. The topic-comment
articulation is therefore prior to the agent--patient asymmetry, not derived
from it.

The distinction can be illustrated immediately with a minimal aspectual pair,
using the verb \ird{kup-} \trsl{buy}:

\pex\label{ex:tc-topic-subj}
\a
\begingl
\gla Marek tak tóme kupujime.//
\glb Marek.\Dir{} \Dem{}.\Prox{}.\Dir{} book.\Ind{} buy-\Cont{}//
\glft \trsl{Marek is buying this book.}
      \lingcontext{nominative domain; agent = absolutive pivot = topic}//
\endgl
\a
\begingl
\gla Tak tóm Marku koupnik.//
\glb \Dem{}.\Prox{}.\Dir{} book.\Dir{} Marek.\Trs{} buy-\Trs{}.\Pf{}//
\glft \trsl{This book, Marek bought.}
      \lingcontext{ergative domain; patient = absolutive pivot = topic}//
\endgl
\xe

\noindent The event participants are constant: Marek is the buyer and the book
is the thing bought. What changes is not the lexical event structure but the
alignment: the nominative-domain continuous places the agent in the direct case
and therefore makes the agent the absolutive pivot and clause topic; the
ergative-domain perfective places the patient in the direct case and makes the
patient the absolutive pivot and clause topic. A notion of `topic' that reduced
to `subject' could not handle this alternation without positing two
grammatically unrelated subject positions. The present analysis handles it as a
single topic mechanism --- the discourse prominence of the absolutive pivot ---
interacting with the aspectual determination of which participant receives the
direct case.

The claim that topic is the logical subject of predication also entails that the
\term{comment} is a semantic predicate. In (\ref{ex:tc-topic-subj}a), the topic
is Marek and the comment is the property of being engaged in the act of buying
this book. In (\ref{ex:tc-topic-subj}b), the topic is the book and the comment
is the property of having been bought by Marek. The predication relation in both
cases is the same type, even though the discourse-prominent participant differs.
This is the sense in which the topic-comment structure is the primary
information-structural dimension of the Iridian clause: it organises the clause
into an individual and a property attributed to it, prior to any further
packaging by focus or modal operators.


\subsection{Absolutive pivot, alignment, and default topicality}
\label{sec:tc-abs-align}

The preceding subsection argued that topic and grammatical subject are
independent notions and that the absolutive pivot, not agenthood per se, is the
default discourse centre. This section formalises the connection between the
absolutive and topichood as a named mapping principle and explains its
interaction with aspect domain.

In neutral finite clauses, there is a systematic alignment between morphological
case marking and discourse structure. The direct-case (absolutive) argument
surfaces clause-initially in Spec,\synhead{TopP} and is interpreted as the
clause topic. This is not a semantic entailment --- the direct-case argument is
not logically required to be discourse-prominent --- but the neutral or unmarked
mapping between clause structure and discourse packaging. The generalisation can
be stated as follows:

\begin{description}[leftmargin=2em, labelwidth=!, labelindent=0em]
\item[\upshape Absolutive--Topic Alignment:] In unmarked finite clauses, the
  \synhead{DP} bearing the direct case is the default discourse topic and is
  merged or interpreted in Spec,\synhead{TopP}.
\end{description}

\noindent This principle is not a semantic law and not an inviolable syntactic
constraint. It is the neutral mapping: deviations from it are possible but
marked, and they trigger the resumptive, contrastive, and focus-related
interpretations discussed in
§§\,\ref{sec:tc-marked-topics}--\ref{sec:tc-contrastive-topic}. The principle is
best understood as the Iridian instantiation of a more general interface
tendency: participants that are structurally prominent (absolutive arguments,
clause-initial position) tend also to be discourse-prominent.

The interaction with aspect is crucial. Aspect choice in Iridian is partly a
pragmatic decision about which participant is to be construed as the discourse
centre: the speaker selects the aspectual domain that places the intended topic
in absolutive position. The same lexical predicate can construe the same event
with different discourse centres:

\pex\label{ex:tc-aspect-topic}
\a
\begingl
\gla Marek travě piaštime.//
\glb Marek.\Dir{} bread.\Ind{} eat-\Prog{}//
\glft \trsl{Marek is eating bread.}
      \lingcontext{nominative domain; agent = direct = topic}//
\endgl
\a
\begingl
\gla Travat Marku piaštnik.//
\glb bread.\Dir{}.\Def{} Marek.\Trs{} eat-\Trs{}.\Pf{}//
\glft \trsl{Marek ate the bread.}
      \lingcontext{ergative domain; patient = direct = topic}//
\endgl
\xe

\noindent The alternation is not mere word-order variation. In
(\ref{ex:tc-aspect-topic}a), the aspectual morphology selects the nominative
domain and the agent surfaces in the direct case. The clause is organised around
Marek as its discourse centre: it is a clause about Marek's ongoing activity. In
(\ref{ex:tc-aspect-topic}b), the aspectual morphology selects the ergative
domain and the patient surfaces in the direct case. The clause is now about the
bread: specifically, about what happened to it. The ergative agent is present as
a core participant but is secondary from a discourse perspective --- it occupies
the transitive case and remains within the comment domain unless specially
promoted (§\,\ref{sec:tc-marked-topics}).

In the imperfective aspects (progressive, continuous, contemplative), agent and
topic therefore typically coincide; in the perfective and retrospective, patient
and topic typically coincide. This is one of the most theoretically
consequential properties of the Iridian grammar and is not merely a descriptive
observation about word order. It reflects the interaction of two independently
motivated systems --- the aspectual specification of which participant receives
the absolutive and the discourse mapping that places the absolutive in topic
position --- and it means that the topic-comment system in Iridian is
inextricably bound up with split-ergative alignment (see
§§\,\ref{sec:th-topic-comment-theory}, \ref{sec:th-wh-interrogation}, and
Chapter~\ref{ch:discourse}).


\subsection{The formal properties of ordinary topics}
\label{sec:tc-ordinary-topics}

Not every DP that might conceivably appear clause-initially is a topic, and not
every type of DP is equally available as a topic. Topics in the unmarked case
have specific formal properties, and distinguishing these properties from those
of other clause-initial material is essential for the analysis of the left
periphery as a whole.

An \term{ordinary topic} in Iridian is a \synhead{DP} that is
discourse-identifiable, referential or kind-denoting, and at least weakly
specific or recoverable in context. These properties are not independently
stipulated features; they follow from the semantic role of the topic as the
logical subject of predication. A predication requires an individual or kind
about which something is being said: predicating a property of an entity that
has not been introduced into, or cannot be recovered from, the discourse would
fail to connect the clause to its context of utterance. The ordinary topic is
therefore inherently discourse-anchored, though the nature and degree of that
anchoring may vary by topic type.

Four principal topic types can be distinguished within the class of ordinary
topics. \emph{Anchored referential \synhead{DP}s} are the paradigm case: DPs
bearing the anchoring suffix \ird{-to}, which signals that the referent is
tracked as a specific individual within the discourse. Anchored referential
topics may be discourse-anaphoric, situationally unique, or associated with a
contextually salient set. The anchoring suffix is by no means required on every
topic (§\,\ref{sec:tc-anchoring-interface}), but anchored topics are the most
prototypically discourse-linked type. \emph{Discourse-recoverable null topics}
arise when the topic referent is established by the prior discourse and
Spec,\synhead{TopP} is phonologically empty; these are treated in
§\,\ref{sec:tc-topic-drop}. \emph{Generic and kind-denoting \synhead{DP}s} may
appear as topics when interpreted as kinds or as generic individuals: the
predication is then a characterising statement about the kind, and the topic is
discourse-stable once introduced. \emph{Discourse-linked specific indefinites}
may appear as topics when the referent is being introduced as particular and
contextually distinguishable; in Iridian, the specifically indefinite particles
\ird{si} and \ird{druh} support this function (§\,\ref{sec:nm-spec}).

The following examples illustrate each type:

\pex\label{ex:tc-topic-types}
\a\label{ex:tc-anch-topic}
\begingl
\gla Marko v zdravime.//
\glb Marek.\Dir{}.\Def{} already sleep-\Prog{}//
\glft \trsl{Marek is already sleeping.} \lingcontext{anchored referential topic}//
\endgl
\a\label{ex:tc-null-topic}
\begingl
\gla By v zdravime?//
\glb \textsc{q} already sleep-\Prog{}//
\glft \trsl{Is he already sleeping?} \lingcontext{null continuing topic; referent established}//
\endgl
\a\label{ex:tc-generic-topic}
\begingl
\gla Vlak drže na hlučnime.//
\glb wolf.\Dir{} night in howl-\Cont{}//
\glft \trsl{Wolves howl at night.} \lingcontext{generic topic; bare plural}//
\endgl
\a\label{ex:tc-spec-indef-topic}
\begingl
\gla druh člověk za dveřmi stojime.//
\glb \Indef{} person.\Dir{} behind door-\Pl{}.\Ind{} stand-\Cont{}//
\glft \trsl{A certain person is standing at the door.}
      \lingcontext{individuated indefinite topic; \ird{druh} marks individuation without anchoring}//
\endgl
\xe

\noindent In (\ref{ex:tc-spec-indef-topic}), the modifier \ird{druh} does not
make the referent discourse-recoverable in the anaphoric sense, but it
individuates the referent as a particular entity with contextual grounds for
its existence. This partial individuation is sufficient to support the
predication relation (§\,\ref{sec:nm-spec}). An unindividuated existential
predication --- the assertion that there is merely some entity, without any
individualising implication --- does not have a topic in this sense; that
construction is treated in §\,\ref{sec:tc-topicless}.


\subsection{Topic position in the clause spine}
\label{sec:tc-clause-spine}

The preceding subsections have presupposed that the topic occupies a designated
structural position. The topic is the \synhead{DP} in Spec,\synhead{TopP},
where \synhead{TopP} is the highest argument-hosting projection in the
articulated left periphery of the Iridian clause. The clause spine, set out in
Table~\ref{tab:th-clause-spine} and Figure~\ref{fig:th-left-periphery}, runs:

\begin{center}
\synhead{SpeechActP} $>$ \synhead{EvidP} $>$ \synhead{TopP} $>$
\synhead{PolP} $>$ \synhead{FocP} $>$ \synhead{IP} $>$
\synhead{PhasP} $>$ \synhead{PredP}
\end{center}

\noindent The topic is in Spec,\synhead{TopP} (projection 8). The comment
corresponds to the domain that \synhead{TopP} scopes over: everything from
\synhead{PolP} downward, including the focus position, the inflectional domain,
and the predicate field. \synhead{TopP} is dominated only by the evidential and
speech-act projections, which are not argumental and do not belong to the
comment.

This architecture has several immediate consequences. Evidential particles such
as \ird{oče} occupy \synhead{EvidP} outside and above the topic, not inside the
comment. The clause-typing particle \ird{by} is in Pol\textsuperscript{0},
below the topic and above the focus position; in polar questions, the topic
may remain in Spec,\synhead{TopP} while \ird{by} takes scope over the
predication within its domain. Focused constituents appear in
Spec,\synhead{FocP}, below the topic and within the comment domain, which
explains why a focused constituent and a topic may co-occur without structural
conflict (§\,\ref{sec:tc-focus-contrast}).

Whether the topic in Spec,\synhead{TopP} arrives there by movement from a
base-generated position within \synhead{IP}, or whether it is base-generated in
that position with its predication role established by a long-distance Agree
relation, is a question deliberately left open here. Both derivational routes
are architecturally consistent with the clause spine
(§\,\ref{sec:th-topic-comment-theory}); the present description is committal
about the surface position of the topic and about the topic-comment
interpretive division, while remaining agnostic about the precise derivational
mechanism pending further evidence from reconstruction and binding asymmetries.


\section{Topic and the information-structural field}\label{sec:tc-isf}

\subsection{Topic versus focus}
\label{sec:tc-focus-contrast}

Topic and focus are distinct operations in the Iridian left periphery and must
not be conflated. The distinction is architectural: \synhead{TopP} is the topic
projection, and \synhead{FocP} --- the position immediately above \synhead{IP}
and below \synhead{PolP} --- is the focus projection. The two positions are not
interchangeable, and the constituents that occupy them have different discourse
functions.

A topic is discourse-given or at least discourse-linkable: it establishes the
individual about which the comment is a predication, and it typically introduces
no new information in itself. A focus, by contrast, is the informationally or
contrastively prominent element within the comment domain: it identifies the
relevant variable in an answer to a question under discussion, or marks the
constituent where the scalar or exhaustive interpretation is concentrated. The
focus is an operator whose scope is internal to the comment; the topic is the
frame within which that comment is interpreted.

The two may co-occur, since topichood and focus are properties of different
projections:

\pex\label{ex:tc-top-foc}
\a\label{ex:tc-top-only}
\begingl
\gla Marek tak tóma obrojime.//
\glb Marek.\Dir{} \Dem{}.\Prox{}.\Dir{} book.\Ind{} read-\Cont{}//
\glft \trsl{Marek is reading this book.}
      \lingcontext{topic only; focus in situ at comment-final position}//
\endgl
\a\label{ex:tc-top-plus-foc}
\begingl
\gla Marko tak tóma obrojime.//
\glb Marek.\Dir{}.\Def{} \Dem{}.\Prox{}.\Dir{} book.\Dir{} read-\Cont{}//
\glft \trsl{As for Marek, it is this book he is reading.}
      \lingcontext{topic \ird{Marko}; \ird{tak tóm} fronted to Spec,\synhead{FocP}}//
\endgl
\xe

\noindent In (\ref{ex:tc-top-only}), Marek is the topic and the comment
identifies his ongoing activity; the focus is in situ at the comment-final
position (§\,\ref{sec:th-focus}). In (\ref{ex:tc-top-plus-foc}), Marek remains
the topic --- the clause is still \emph{about} Marek --- but \ird{tak tóm} has
been fronted to Spec,\synhead{FocP}, marking it as the focus: specifically this
book is what Marek is reading. The topic is the discourse anchor; the focus is
the marked operator within the comment.

Several further contrasts follow from this architecture. A topic must be
discourse-linkable; a focus need not be. A topic cannot by itself be the locus
of new, unpredicted information; a focus can. A topic takes scope over the
entire comment domain; a focus takes scope over the predicate field within that
domain. These distinctions are maintained throughout the grammar, and material
at the left periphery should always be identified as topic, focus, adverbial, or
clause-typing particle before any structural claim is made about it.


\subsection{Topic versus sentence adverbials and peripheral particles}
\label{sec:tc-adverbials}

One of the most practically important distinctions in the left periphery is
between true topics and other clause-peripheral material. It is tempting to
treat every clause-initial constituent as a topic simply because it precedes the
predicate. This is incorrect. A true topic is an argumental constituent that
stands in a predication relation with the comment: the comment predicates a
property of the topic, and the topic is the semantic individual about which that
predication holds. Sentence adverbials, evidential particles, and clause-typing
operators have entirely different grammatical and discourse functions and should
not be analysed as topics merely because they appear to the left of the comment
domain.

A \term{sentence adverbial} --- a temporal, locative, or modal adjunct that
scopes over the proposition --- is not a constituent of the argument structure
and is not the individual about which the comment predicates something.
\ird{Bych} \trsl{yesterday}, for instance, provides a temporal frame for the
event; the clause does not predicate a property \emph{of yesterday}. Similarly,
\term{evidential particles} such as \ird{oče} occupy \synhead{EvidP}, which is
above \synhead{TopP}: they contribute an epistemic or evidential stance on the
proposition as a whole. These material types are licensed by their own
projections in the left periphery and must be kept analytically separate from
topical material.

The following examples illustrate three distinct clause-initial configurations:

\pex\label{ex:tc-adv-contrast}
\a\label{ex:tc-adv-true-top}
\begingl
\gla Marek by v zdravime?//
\glb Marek.\Dir{} \textsc{q} already sleep-\Prog{}//
\glft \trsl{Is Marek sleeping yet?}
      \lingcontext{true topic; \ird{Marek} in Spec,\synhead{TopP}}//
\endgl
\a\label{ex:tc-adv-adv}
\begingl
\gla Bych by Marek v zdravime?//
\glb yesterday \textsc{q} Marek.\Dir{} already sleep-\Prog{}//
\glft \trsl{Was Marek already asleep yesterday?}
      \lingcontext{temporal adverbial \ird{bych} above \synhead{PolP}; \ird{Marek} is still the topic}//
\endgl
\a\label{ex:tc-adv-evid}
\begingl
\gla Oče Marko v zdravime.//
\glb \textsc{evid} Marek.\Dir{}.\Def{} already sleep-\Prog{}//
\glft \trsl{Apparently Marek is already sleeping.}
      \lingcontext{evidential \ird{oče} in Spec,\synhead{EvidP} above topic; \ird{Marko} is the topic}//
\endgl
\xe

\noindent In (\ref{ex:tc-adv-true-top}), \ird{Marek} is the topic: the clause
predicates a state of him. In (\ref{ex:tc-adv-adv}), \ird{bych} is a temporal
adverbial adjoined above \synhead{PolP}; it provides the temporal frame. The
clause still has a topic --- \ird{Marek} --- which follows the adverbial in the
linear string. In (\ref{ex:tc-adv-evid}), \ird{oče} is the evidential particle
in Spec,\synhead{EvidP}, scoping over the whole proposition including the topic.
The topic is still \ird{Marko}.

The practical consequence is significant: the left edge of an Iridian clause may
contain a stack of material from distinct projections --- a speech-act particle,
then an evidential, then a topic, then a clause-typing particle. Only the
constituent in Spec,\synhead{TopP} is the logical subject of predication. The
remainder should be analysed in their own designated positions.


\section{Topic in discourse: drop, absence, and multiplicity}\label{sec:tc-discourse}

\subsection{Multiple topics and the hierarchy of the left periphery}
\label{sec:tc-multiple-topics}

The architecture of \synhead{TopP} allows for the possibility that more than one
topic-level projection is active in a clause. Ordinary Iridian clauses have a
single topic in Spec,\synhead{TopP}, and iterated topic structure is not the
default. But the grammar appears to permit a second topic projection in specific
discourse environments, particularly where a \term{frame-setting topic}
establishes a domain of quantification or a contrastive partition, and a lower
\term{clause topic} identifies the individual about whom the comment predicates
something within that frame.

A frame-setting topic is a constituent that restricts the discourse domain in
which the comment is interpreted. It does not name the individual that the
comment is directly about; rather, it sets up a restricted context within which
the clause topic is identified and the comment evaluated. Frame-setting topics
are often formally marked by prosodic independence and, in some contexts, by a
resumptive element within the lower clause that co-refers with the frame:

\pex\label{ex:tc-frame-top}
\a\label{ex:tc-frame-town}
\begingl
\gla Tak dvor Marko v něm bydlime.//
\glb \Dem{}.\Prox{}.\Dir{} town.\Dir{} Marek.\Dir{}.\Def{} already in.it live-\Cont{}//
\glft \trsl{As for this town, Marek already lives in it.}
      \lingcontext{frame topic + clause topic; resumptive \ird{něm} co-refers with frame topic}//
\endgl
\a\label{ex:tc-frame-marek}
\begingl
\gla Marko tak tóm bych čitanik.//
\glb Marek.\Dir{}.\Def{} \Dem{}.\Prox{}.\Dir{} book.\Dir{} yesterday read-\Trs{}.\Pf{}//
\glft \trsl{As for Marek, this book he read yesterday.}
      \lingcontext{frame topic + ergative-domain absolutive topic}//
\endgl
\xe

\noindent In (\ref{ex:tc-frame-town}), \ird{tak dvor} \trsl{this town} is the
frame topic and \ird{Marko} is the clause topic; the resumptive pronoun
\ird{něm} inside the comment co-refers with the frame topic. In
(\ref{ex:tc-frame-marek}), \ird{Marko} is the frame topic and \ird{tak tóm}
is the patient-absolutive topic of the ergative clause, within the frame of
what Marek has done.

The distributional limits of iterated \synhead{TopP} await more detailed
investigation. The structural possibility is licensed by the cartographic
architecture: since \synhead{TopP} is a full projection with a specifier and a
complement, it may in principle be embedded within a higher \synhead{TopP} whose
specifier hosts the frame-setting topic. This analysis is consistent with
evidence from topic continuity and topic shift in extended discourse
(Chapter~\ref{ch:discourse}).


\subsection{Topic drop and null topics}
\label{sec:tc-topic-drop}

Topic drop --- the absence of an overt topic \synhead{DP} when the referent is
established by the prior discourse --- is the default strategy for topic
continuation in Iridian. It is not a separate construction or the result of a
dedicated deletion rule; it is the ordinary null spellout of Spec,\synhead{TopP}
under discourse-conditioned recoverability. The theoretical point is important:
topic drop does not reduce the clause to a structure without a topic projection.
Rather, Spec,\synhead{TopP} is syntactically present and semantically
interpreted --- its referential value is supplied by the discourse context ---
but it receives no phonological exponent.

The licensing condition is discourse recoverability. When the most recently
topicalised referent has not shifted and remains available as the intended
individual, Spec,\synhead{TopP} may be null. When the topic shifts, an overt
\synhead{DP} is typically required to signal the new discourse centre, though a
fronted pronoun with contrastive or resumptive force is also available. The
interaction with Zone~\textsc{i} particles follows directly: since those
particles occupy positions post-topic in the left periphery
(§\,\ref{sec:th-topic}), they appear at the left edge of the string when the
topic is null, without any separate movement being involved.

The following discourse sequence illustrates the pattern:

\pex\label{ex:tc-topdrop}
\a\label{ex:tc-topdrop-overt}
\begingl
\gla Marko domů chodime.//
\glb Marek.\Dir{}.\Def{} home go-\Cont{}//
\glft \trsl{Marek is going home.} \lingcontext{overt topic; referent introduced}//
\endgl
\a\label{ex:tc-topdrop-null1}
\begingl
\gla by v domě jest?//
\glb \textsc{q} already at.home be//
\glft \trsl{Is he home yet?}
      \lingcontext{null topic; referent continues from (\ref{ex:tc-topdrop-overt})}//
\endgl
\a\label{ex:tc-topdrop-null2}
\begingl
\gla ne, ještě cestujime.//
\glb no still travel-\Cont{}//
\glft \trsl{No, he is still travelling.}
      \lingcontext{null topic continues; negation particle \ird{ne}}//
\endgl
\xe

\noindent The null topic in
(\ref{ex:tc-topdrop-null1}--\ref{ex:tc-topdrop-null2}) is semantically
interpreted as co-referential with \ird{Marko} from (\ref{ex:tc-topdrop-overt}).
No referential shift has occurred, and no overt \synhead{DP} is needed. The
question particle \ird{by} and the phasal particle \ird{v} appear at the left
edge of the string because their structural positions are post-topic, not
because any movement has targeted them.

Topic drop should be distinguished from radical pro-drop in languages that lack
agreement morphology altogether. In Iridian, there is no verbal person
agreement; the recoverability of null topics therefore depends entirely on
discourse tracking and the prior absolutive pivot. The topic most naturally
continued by zero anaphora is the most recent absolutive argument, not simply
the most salient animate participant (§\,\ref{sec:th-topic};
Chapter~\ref{ch:discourse}). This makes the absolutive status of the referent
--- not its agentive or grammatical-subject status --- the crucial factor in
licensing null continuation.


\subsection{Apparent and genuinely topicless clauses}
\label{sec:tc-topicless}

Not every clause without an overt topic \synhead{DP} is structurally topicless.
The preceding section showed that null topics are the ordinary form of topic
continuation under discourse recoverability. A clause that lacks an overt topic
may therefore be either (a) a clause whose topic is syntactically present but
phonologically null, or (b) a clause that genuinely lacks a topic --- one in
which Spec,\synhead{TopP} is not projected, or is projected without an argument.
These two types must be distinguished because they have different discourse
properties and different structural analyses.

\term{Apparent topicless clauses} are clauses in which Spec,\synhead{TopP} is
phonologically null but semantically filled. The null topic has a referential
value recoverable from the discourse; it supports further topic continuation by
zero anaphora; and it behaves structurally like any overt topic with respect to
aspect selection and the ordering of comment-internal material.
Examples~(\ref{ex:tc-topdrop-null1}--\ref{ex:tc-topdrop-null2}) above are
apparent topicless clauses of this type.

\term{Genuinely topicless clauses} are clauses in which the predication does not
have a pre-established discourse individual as its logical subject. These
include existential and presentational constructions --- clauses whose function
is to introduce an entity or situation into the discourse rather than to
predicate a property of an already established individual. In such clauses there
is no prior referent to anchor a topic, and the clause as a whole asserts the
existence or occurrence of something. The existential construction with
\ird{jec-} is the primary example (§\,\ref{sec:existential-constructions}):

\pex\label{ex:tc-topicless}
\a\label{ex:tc-apparent-topicless}
\begingl
\gla by koupnik?//
\glb \textsc{q} buy-\Trs{}.\Pf{}//
\glft \trsl{Did he buy it?}
      \lingcontext{apparent topicless; null topic and null object, both recoverable}//
\endgl
\a\label{ex:tc-pres-topicless}
\begingl
\gla jec druh člověk za dveřmi.//
\glb \Exst{} \Indef{} person.\Dir{} behind door-\Pl{}.\Ind{}//
\glft \trsl{There is a certain person at the door.}
      \lingcontext{genuinely topicless presentational; \ird{druh} individuates; entity introduced into discourse}//
\endgl
\a\label{ex:tc-exist-topicless}
\begingl
\gla jec hluk v dvoře.//
\glb \Exst{} noise.\Dir{} in yard.\Ind{}//
\glft \trsl{There is noise in the yard.}
      \lingcontext{genuinely topicless existential-eventive; situation introduced}//
\endgl
\xe

\noindent In (\ref{ex:tc-apparent-topicless}), both the topic \synhead{DP} and
the object are null but both are recoverable: the buyer and the thing bought
have already been established in the discourse. Spec,\synhead{TopP} is present
but phonologically empty. In
(\ref{ex:tc-pres-topicless}--\ref{ex:tc-exist-topicless}), by contrast, the
\ird{jec-} clauses introduce new referents into the discourse. There is no prior
individual about whom the clause predicates a property; the person at the door
and the noise in the yard are the entities being asserted to exist, not the
continuation of an established topic. These clauses do not project an
argumentally filled Spec,\synhead{TopP}.

Genuinely topicless clauses are not defective or marginal; they are a
grammatically well-defined type in Iridian, closely associated with the
\ird{jec-}/\ird{zad-} existential paradigm, with certain environmental
predicates, and with presentational discourse moves that introduce entities for
the first time. The overlap between the existential construction and the
genuinely topicless clause is one of the places where the topic-comment system
and the existential chapter (§\,\ref{sec:existential-constructions}) mutually
support one another rather than remaining separate descriptive domains.


\section{Non-default topics: marked, promoted, and contrastive}\label{sec:tc-nondefault}

\subsection{Marked topics: non-direct, resumptive, and promoted arguments}
\label{sec:tc-marked-topics}

The Absolutive--Topic Alignment principle (§\,\ref{sec:tc-abs-align}) identifies
the direct-case pivot as the default topic. All departures from this default are
marked: they signal that the speaker has a specific discourse motivation for
presenting a non-absolutive participant as the clause's discourse centre, rather
than accepting the mapping that aspect domain provides. Marked topics are
therefore not free variants of ordinary topics; they carry additional
interpretive weight and are associated with specific discourse functions.

The nominal chapter establishes
(§§\,\ref{sec:nm-topic-interface}--\ref{sec:nm-topic-oblique}) that anchoring on
transitive-case and indirect-case \synhead{DP}s is always more plainly
expressive of contrastive, resumptive, or topic-shift function than anchoring on
direct-case \synhead{DP}s. The reason is structural: a transitive-case agent is
not the absolutive pivot and is not in Spec,\synhead{TopP} under the neutral
mapping; if such an agent bears the anchoring suffix, the anchoring performs an
additional function --- signalling the agent's individuation or promotion ---
beyond the baseline discourse-tracking role that anchoring has on direct-case
\synhead{DP}s.

In the ergative domain, the neutral clause has the patient as absolutive topic.
A speaker who wishes instead to make the agent the discourse centre may either
select a nominative-domain aspect (which makes the agent the absolutive pivot by
aspect choice), or produce a marked clause in which the agent is specially
fronted or anchored. The latter creates a marked agent topic. Compare:

\pex\label{ex:tc-marked-top}
\a\label{ex:tc-neutral-top}
\begingl
\gla Tak tóm Marku koupnik.//
\glb \Dem{}.\Prox{}.\Dir{} book.\Dir{} Marek.\Trs{} buy-\Trs{}.\Pf{}//
\glft \trsl{This book, Marek bought.}
      \lingcontext{neutral ergative clause; patient = absolutive = topic}//
\endgl
\a\label{ex:tc-marked-agent-top}
\begingl
\gla Markuto tak tóm koupnik.//
\glb Marek-\Def{}.\Trs{} \Dem{}.\Prox{}.\Dir{} book.\Dir{} buy-\Trs{}.\Pf{}//
\glft \trsl{As for Marek (himself), he bought this book.}
      \lingcontext{marked: anchored agent in transitive case promoted to topic field}//
\endgl
\xe

\noindent In (\ref{ex:tc-neutral-top}), the book is the absolutive pivot and the
clause topic; Marek occupies the transitive case in the comment domain. In
(\ref{ex:tc-marked-agent-top}), the anchored agent \ird{Markuto} has been
promoted to or interpreted in the topic field, while the book retains the direct
case. This is not the ordinary ergative mapping: the clause signals agent
resumption, contrast, or local topic shift. The interpretive effect is typically
one of emphasis or correction --- it was Marek, not someone else, who bought
this book --- or a switch of discourse tracking onto Marek as the continuing
topic.

A directly comparable construction involving an indirect-case argument is also
available when a recipient or beneficiary is being individuated as the most
discourse-prominent participant. Such clauses are marked by the same criteria
and are developed further in the discourse chapter (Chapter~\ref{ch:discourse}).


\subsection{Contrastive topic}
\label{sec:tc-contrastive-topic}

Contrastive topic is a distinct phenomenon that should not be reduced to
ordinary topic plus incidental emphasis, nor to focus. A \term{contrastive
topic} is a topic bearing a contrastive feature: it establishes the clause as
being about an individual X while simultaneously invoking a salient set of
alternatives to X, implying that the comment may not hold of all members of that
set. The structure is one of partial predications: the clause provides
information about X, and the contrastive framing signals that the corresponding
predication about the alternatives to X remains open or is differently answered.

This differs from focus in a fundamental respect. In a focused clause, the focus
constituent exhaustively identifies the relevant value within the comment (it is
\emph{this book} that Marek is reading, not some other): the identification is
within the predicate domain. In a contrastive topic construction, the
contrastive element is the subject of the predication, not the identified value
within the predicate: it is \emph{Marek} (among others) about whom this comment
holds. The alternatives are not alternatives to the complement of the predicate
but alternatives to the topic individual.

\pex\label{ex:tc-ctop}
\a\label{ex:tc-ctop-ordinary}
\begingl
\gla Marko tak tóme čitajime.//
\glb Marek.\Dir{}.\Def{} \Dem{}.\Prox{}.\Dir{} book.\Ind{} read-\Cont{}//
\glft \trsl{Marek is reading this book.}
      \lingcontext{ordinary topic; no contrastive alternatives invoked}//
\endgl
\a\label{ex:tc-ctop-contr}
\begingl
\gla Marko tak tóme čitajime, ale Janko ne.//
\glb Marek.\Dir{}.\Def{} \Dem{}.\Prox{}.\Dir{} book.\Ind{} read-\Cont{} but Janek.\Dir{}.\Def{} \Neg{}//
\glft \trsl{Marek is reading this book; Janek is not.}
      \lingcontext{contrastive topic \ird{Marko}; Marek contrasted with Janek}//
\endgl
\a\label{ex:tc-ctop-obj}
\begingl
\gla Tak tóm Marku koupnik, ale si tuto list ne.//
\glb \Dem{}.\Prox{}.\Dir{} book.\Dir{} Marek.\Trs{} buy-\Trs{}.\Pf{} but \Sanch{} \Dem{}.\Med{}.\Dir{} letter.\Dir{} \Neg{}//
\glft \trsl{This book, Marek bought; a certain particular letter, he did not.}
      \lingcontext{contrastive topic in ergative-domain clause; \ird{si} speaker-anchors the letter}//
\endgl
\xe

\noindent In (\ref{ex:tc-ctop-ordinary}), Marek is an ordinary topic: the clause
predicates reading of him, but no alternatives are invoked. In
(\ref{ex:tc-ctop-contr}), the contrastive feature on \ird{Marko} signals that
the addressee should consider the contrast with Janek: it is Marek (as opposed
to Janek) about whom this reading predication holds. In (\ref{ex:tc-ctop-obj}),
the contrastive topic is \ird{tak tóm} in an ergative-domain clause: this book
(as opposed to a particular letter) is the thing Marek bought. The contrastive
topic here is the patient-absolutive, the ordinary topic of an ergative clause,
bearing an additional contrastive feature.

In the phonological component of the grammar, contrastive topics typically
receive a distinct prosodic contour, the analysis of which is deferred to the
prosody chapter. What the present section establishes is that contrastive topics
are syntactically and discourse-structurally distinct from both ordinary topics
and from focus, and that Spec,\synhead{TopP} is the structural locus of both
ordinary and contrastive topics, with the contrastive feature as the
distinguishing property.


\section{The topic--anchoring interface}
\label{sec:tc-anchoring-interface}

The interaction between clause-level topichood and the nominal anchoring system
is treated in detail in Chapter~\ref{ch:nouns}
(§\,\ref{sec:nm-topic-interface}), where the \term{Topicality--Anchoring
Harmony} preference is stated: a \synhead{DP} bearing the feature
[\textsc{+top}] in Spec,\synhead{TopP} is preferentially realised with the
anchoring suffix \ird{-to} in its nominal functional hierarchy. The present
section draws out the descriptive consequence for the exposition of
topic-comment structure.

The preference rests on a straightforward functional logic. Topics are
discourse-salient, discourse-trackable, and identified individuals. These are
precisely the conditions that the anchoring suffix encodes: \ird{-to} signals
that the referent is an anchored discourse token, tracked within the common
ground (§\,\ref{sec:nm-def-semantics}). The strongest co-occurrence of topic and
anchoring is therefore expected with ordinary referential topics. The preference
is weaker with generic and bare nominal topics, where the relevant entity is a
kind rather than a particular, and absent with null topics, which are by
definition recoverable without overt morphology.

A minimal contrast makes the distinction visible:

\pex\label{ex:tc-anch-contrast}
\a\label{ex:tc-anch-overt}
\begingl
\gla Marko v zdravime.//
\glb Marek.\Dir{}.\Def{} already sleep-\Prog{}//
\glft \trsl{Marek is already sleeping.}
      \lingcontext{anchored topic; explicit discourse tracking}//
\endgl
\a\label{ex:tc-anch-bare}
\begingl
\gla Marek v zdravime.//
\glb Marek.\Dir{} already sleep-\Prog{}//
\glft \trsl{Marek is already sleeping.}
      \lingcontext{unanchored topic; less explicit tracking}//
\endgl
\xe

\noindent The two clauses are truth-conditionally equivalent. The difference is
discourse-structural: (\ref{ex:tc-anch-overt}) presents Marek as a trackable
individual whose identity is being flagged in the common ground, whereas
(\ref{ex:tc-anch-bare}) presents him as the clause topic without that additional
marking. The anchoring suffix does not change the propositional content; it
adjusts the degree to which the topic is explicitly registered as a discourse
token. The Topicality--Anchoring Harmony preference is violable: bare nominal
topics, generic topics, and contrastive topics may all appear without overt
anchoring. The preference captures a statistical correlation with a structural
explanation, not a categorical requirement.


\section{Summary}
\label{sec:tc-summary}

The analysis of topic-comment structure developed in this section can be
summarised in the following claims:

\begin{enumerate}[wide]
\item The Iridian clause is organised primarily by a \term{topic-comment
  bipartition}. The topic is the logical subject of predication; the comment is
  the semantic predicate attributed to the topic. This articulation is
  independent of grammatical subjecthood.

\item The default topic of a neutral finite clause is the direct-case absolutive
  pivot occupying Spec,\synhead{TopP}: the \term{Absolutive--Topic Alignment}.
  In the imperfective aspects the absolutive is the agent; in the perfective and
  retrospective it is the patient. The topic-comment system is therefore
  inextricably tied to split-ergative alignment.

\item Ordinary topics are discourse-identifiable, referential or kind-denoting,
  and at least weakly specific or recoverable. The four principal types are
  anchored referential \synhead{DP}s, null continuing topics, generics, and
  discourse-linked specific indefinites.

\item \synhead{TopP} is the highest argument-hosting projection in the clause
  spine. The comment corresponds to the domain below \synhead{TopP}. Evidential
  and speech-act projections are above \synhead{TopP} and do not belong to the
  comment.

\item Topic and focus are distinct left-peripheral operations: the topic is the
  discourse anchor; the focus is the informationally or contrastively prominent
  element within the comment. The two may co-occur in their respective
  projections.

\item Sentence adverbials, evidential particles, and clause-typing particles are
  not topics; they occupy their own designated projections.

\item \synhead{TopP} may be iterated in frame-setting and resumptive contexts,
  yielding a frame-setting topic and a lower clause topic.

\item Topic drop is the ordinary null spellout of Spec,\synhead{TopP} under
  discourse recoverability, not the absence of topic structure.

\item A clause without an overt topic may be an apparent topicless clause (null
  topic present, recoverable from discourse) or a genuinely topicless clause (no
  topic projected). Genuinely topicless clauses are associated primarily with
  existential and presentational predicates.

\item Marked topics --- anchored non-direct \synhead{DP}s, promoted agents,
  individuated recipients --- are departures from the Absolutive--Topic
  Alignment signalling resumption, contrast, or switch-tracking.

\item Contrastive topics bear a contrastive feature, invoking alternatives to
  the topic individual; they differ from both ordinary topics and from focus.

\item Topics are preferentially anchored because anchoring and topicality are
  functionally correlated; this preference is strongest with referential topics
  and absent with null topics.
\end{enumerate}

